% Included by the SOS operating manual. All links resolve inside the manual.
\soschapter{dictionary}{Terminology dictionary}{Reference}{Key terminology and notation}
A name, a cell coordinate and a configuration address answer different lookup questions. The entries below explain each term in the operational interface and the symbols used in the original equations.

\noindent\textsf{\hyperlink{g-a}{A--C}\quad\hyperlink{g-d}{D--H}\quad
\hyperlink{g-i}{I--N}\quad\hyperlink{g-o}{O--Q}\quad
\hyperlink{g-r}{R--S}\quad\hyperlink{g-t}{T--Z}\quad
\hyperlink{g-notation}{Mathematical signs}}

\newcommand{\sosdictgroup}[2]{\par\Needspace{5\baselineskip}\manualtarget{#1}\subsection*{#2}}

\sosdictgroup{g-a}{A--C}
\sosterm{g-addition}{Addition, $\oplus$}{Combine two configurations by adding matching coordinates. With $x=(1,0)$ and $y=(0,1)$, $x\oplus y=(1,1)$. Zero is the identity. Addition admits cancellation, meaning $x\oplus z=y\oplus z$ implies $x=y$.}
\sosterm{g-admissible}{Admissible}{Allowed by the selected profile. In the natural profile every coordinate is a nonnegative integer. A restricted profile can impose additional conditions on whole configurations. An operation on admitted inputs may still produce an excluded result. See \hyperlink{g-restricted}{restricted domain}.}
\sosterm{g-lattice}{Affine lattice}{A set of integer vectors consisting of a particular solution plus every integer combination of some direction vectors. For signed neutral--toggle division, $(0,1)+t(1,-1)$ for $t\in\mathbb Z$ describes an entire quotient set.}
\sosterm{g-alphabet}{Alphabet, $A$}{The finite, nonempty ordered list $(o_1,\ldots,o_d)$ of primitive types. Its order fixes which type a coordinate counts and breaks ties in the canonical ranking. The size $d$ counts types, including types whose current multiplicity is zero.}
\sosterm{g-array}{Array, axis and cell}{An axis is one coordinate position in an array index. A cell is the value at one permitted index tuple. The registry array holds records. The array inside a record holds configurations. See \hyperlink{g-shape}{shape and extent}.}
\sosterm{g-associative}{Associativity and commutativity}{Associativity means $(x\star y)\star z=x\star(y\star z)$. The predicate $\mathsf A_T$ tests equality for every triple of basis configurations and every output coordinate. Commutativity means $x\star y=y\star x$. The predicate $\mathsf C_T$ tests $T_{ij}^{k}=T_{ji}^{k}$ for every $i,j,k$. Bilinearity extends these basis tests to all configurations. \linesref{T}{305}{346}.}
\sosterm{g-atomic}{Atomic registration}{Publishing the name index and the corresponding occupied grid cell as one state change. A reader of a snapshot can see either complete registry state. A failed registration leaves the preceding state available.}
\sosterm{g-basis}{Basis configuration, $e_i$}{The vector containing one occurrence at coordinate $i$ and zero elsewhere. Every natural configuration has the unique decomposition $\bigoplus_i x_i e_i$. The expression $n x$ means the sum of $n$ copies of $x$, with zero copies giving the empty configuration.}
\sosterm{g-bijection}{Bijection}{A function pairing each input with one output so that every output has exactly one input. Its inverse reverses the pairing. \emph{Injective} means distinct inputs have distinct outputs. \emph{Surjective} means every member of the declared output domain is reached.}
\sosterm{g-bilinearity}{Bilinearity and distributivity}{Composition distributes over a sum in either operand. Thus every term $x_i y_j$ contributes $x_i y_j$ copies of $e_i\star e_j$ to the product. Both products with zero equal zero, the absorbing element for multiplication.}
\sosterm{g-byte}{Bit, byte and big-endian order}{A bit is zero or one. A byte has eight bits and represents a number from $0$ to $255$. In big-endian order the first byte gives the coefficient of the greatest power of $256$. For example, the bytes \texttt{01 00}, shown in hexadecimal (base sixteen), represent $256$.}
\sosterm{g-rank}{Canonical rank and unranking, $\rho,\delta$}{The rank $\rho(x)$ counts configurations before $x$ in the declared order. The natural profile orders by size and then lexicographically. The inverse $\delta=\rho^{-1}$ recovers the vector. In the rank formula $p_i(x)=\sum_{h<i}x_h$ counts earlier coordinates. To unrank an address, find its size shell and select the vector at its position within that shell. \linesref{T}{125}{185}.}
\sosterm{g-cardinality}{Cardinality, $\nu$}{For a natural configuration $x$, $\nu(x)=\sum_i x_i$ counts occurrences. For a set $U$, $|U|$ counts members. Thus $|\{(1,1)\}|=1$ while $\nu((1,1))=2$. On addresses, $\nu_I=\nu\circ\delta$. Product sizes multiply for all $x,y$ exactly when every primitive pair yields one occurrence.}
\sosterm{g-codec}{Codec, $E$}{A matched encoder and decoder. For SOS text, $E(s)$ is the unsigned integer formed from the single prefix byte \texttt{01} followed by the UTF-8 bytes of $s$. Decode the shortest big-endian representation, check its prefix, then decode the remaining bytes. The empty text has code $1$. Record encoding applies $E$ after framing its fields. See \hyperlink{impl-1}{Implementation~1}.}
\sosterm{g-completion}{Completion and learning flag}{The caller supplies a Boolean $c$, either $0$ or $1$. The adopted $\operatorname{status\_ping}(c)=1-c$ yields $1$ to request further detail and $0$ when complete. Draft lines 15 and 19 both test $1$, leaving $c=0$ unresolved in the original. \linesref{D}{13}{25}.}
\sosterm{g-tensor}{Composition tensor, $T$, and product, $\star$}{An entry $T_{ij}^{k}$ counts occurrences of type $k$ produced by one ordered primitive pair $(o_i,o_j)$. The vector $\mu_T(o_i,o_j)$ has these entries as its coordinates. For configurations, $(x\star y)_k=\sum_{i,j}x_i y_j T_{ij}^{k}$. The entire tensor belongs to the profile.}
\sosterm{g-configuration}{Configuration, $\mathcal C$}{A multiset of typed occurrences represented by a multiplicity vector. In the natural case $\mathcal C=\mathbb N_0^d$. Two vectors are equal exactly when corresponding counts agree. The zero vector $\mathbf 0$ has no occurrences. The signed extension uses integer coefficients.}
\sosterm{g-congruence}{Congruence and semantic equivalence}{Write $x\sim_\sigma y$ when $\sigma(x)=\sigma(y)$. This equivalence is a congruence for an operation when replacing arguments by equivalent ones gives equivalent results. Defining a partial operation on equivalence classes also requires compatible definedness. Cardinality can hide whether a particular type can be removed. \linesref{T}{601}{635}.}
\sosterm{g-order}{Coordinatewise order, $\preccurlyeq$}{The relation $y\preccurlyeq x$ means $y_i\leq x_i$ for every coordinate $i$. It checks availability of every requested type for exact removal. For example, neither $(1,0)$ nor $(0,1)$ lies below the other.}

\sosdictgroup{g-d}{D--H}
\sosterm{g-determinant}{Determinant, $\det M$}{A signed number associated with a square matrix. For a $1$-by-$1$ matrix take its only entry. For larger matrices expand along row one, summing $(-1)^{1+j}M_{1j}$ times the determinant after deleting row one and column $j$. A nonzero determinant makes $Mv=0$ have only the zero rational solution. For a fixed natural divisor and side this is equivalent to at most one quotient for every dividend. A zero determinant permits a collision for some dividend. \linesref{T}{423}{428}.}
\sosterm{g-diagonalization}{Diophantine elimination and pivots}{A Diophantine equation asks for integer solutions. In the signed solver, a pivot is the nonzero entry used to simplify a row. The algorithm swaps columns and subtracts integer multiples of one column from another, preserving an invertible integer coordinate change (a unimodular map). When the row has a single nonzero coefficient, that pivot must divide the corresponding target. Fix its coordinate and solve the remaining rows recursively. See \hyperlink{impl-4}{Implementation~4}.}
\sosterm{g-division}{Division}{Exact multiplicative reconstruction. Declare which side holds the known factor and return all admitted solutions of that equation. See \hyperlink{g-quotient}{quotient set} and \hyperlink{g-matrix}{multiplication matrix}.}
\sosterm{g-domain}{Domain, image and closure}{The domain $\operatorname{dom} f$ contains inputs where $f$ has a value. The image $\operatorname{im} f$ consists of its attained values. A total function is defined for every allowed input. A partial one may be undefined for some. A set is closed under an operation when applying it to members always yields another member.}
\sosterm{g-errors}{Errors and empty results}{\textsf{InvalidInput} rejects malformed arguments, invalid encodings, or numbers outside the declared input domain. \textsf{UnknownName} means a valid name is absent. On registration, \textsf{DuplicateName} means the name already exists and \textsf{OccupiedPosition} means the requested position is taken. Within the algebra, $\mathsf{Fail}$ signifies an undefined scalar operation. $\mathsf{Sol}(\varnothing)$ reports a valid division with no solutions.}
\sosterm{g-evaluator}{Evaluator and outcome tags}{The set of operation symbols $\Omega$ contains $\Omega_{\rm alg}=\{\oplus,\star,\ominus\}$ and the two divisions. $\operatorname{Eval}(\omega,a,b)$ returns a tagged outcome in $\mathsf{Out}_I$. $\mathsf{Val}(a)$ carries one address, $\mathsf{Sol}(U)$ carries a set of addresses, and $\mathsf{Fail}$ records an undefined scalar result. Thus $\mathsf{Val}(0)$ and $\mathsf{Sol}(\{0\})$ have distinct types. \linesref{T}{481}{534}.}
\sosterm{g-fiber}{Fiber}{The complete collection of inputs sent to a given output by a function. Fixing $y$ gives a function $q\mapsto q\star y$. Its fiber over $x$ is $Q_{\rm R}(x,y)$. A fiber can be empty, finite or infinite.}
\sosterm{g-monoid}{Finite monoid and monoid algebra}{A monoid is a set with associative multiplication and a two-sided identity. In $\mathsf M_{m,h}(T)$ the rule $m(i,j)$ chooses the unique output type of each primitive pair, is associative, and has identity $h$. Hence $T_{ij}^{k}=\mathbf1_{\{m(i,j)=k\}}$. Write $M=([d],m,h)$. Its monoid algebra $\mathbb N_0[M]$ uses natural coefficients and the product induced by $m$. The algebra $\mathbb Z[M]$ permits signed coefficients. The unit is $e_h$. \linesref{T}{673}{723}.}
\sosterm{g-framing}{Framing and row-major order}{One field is its canonical decimal UTF-8 byte length, then \texttt{\#}, then exactly that many bytes. For example, \texttt{3\#a\#b} encodes the field \texttt{a\#b}. Record fields are name, profile, axis count, extents and cell coefficients. Visit cells in row-major order, varying the last axis fastest. The shape and alphabet give the number of coefficient fields.}
\sosterm{g-free}{Free commutative monoid, $\operatorname{FreeCMon}(A)$}{The multiset domain under addition and zero. Every map from the alphabet to a commutative monoid extends in exactly one addition- and identity-preserving way to the configurations. Thus counts completely determine each sum of primitives.}
\sosterm{g-free-coordinate}{Free coordinate}{A component of a solution vector that can vary independently within its declared number domain. A zero column in a nonnegative multiplication matrix gives a free natural coordinate whenever a solution exists.}
\sosterm{g-gr}{Group completion, $\operatorname{Gr}$}{Pairs of natural configurations represent formal differences. Pairs $(x,y)$ and $(x',y')$ represent the same difference exactly when $x\oplus y'=x'\oplus y$. Taking one class for each equivalence gives $\mathbb Z^d$ under addition, a group because every vector has an additive inverse. \linesref{T}{815}{841}.}
\sosterm{g-homomorphism}{Homomorphism}{A map preserving the stated operations and identities. For the semantic predicate $\mathsf H_\sigma$, see \hyperlink{g-semantics}{semantic interpretation}. A unital homomorphism also preserves the designated multiplicative identity.}

\sosdictgroup{g-i}{I--N}
\sosterm{g-address}{Identifier and sequential indexing, $I$}{An identifier is the address of one configuration within a namespace. Sequential identifiers use successive labels. In the natural profile $I=\mathbb N_0$ and the canonical order determines which configuration receives each label. A registry assigning the next free label on arrival can be sequential while its labels depend on history.}
\sosterm{g-identity}{Identity, zero and unit}{An identity leaves the other operand unchanged on both sides. The additive identity is $\mathbf0$, with address $0_I=\rho(\mathbf0)$. $\mathcal U_T$ is the set of two-sided multiplicative identities. Checking $u\star e_i=e_i\star u=e_i$ for every $i$ suffices. There is at most one $u$. If it exists its address is $1_I=\rho(u)$, which is $2$ in the neutral--toggle example.}
\sosterm{g-kernel}{Kernel}{For a matrix $M$ on a specified coefficient domain, its zero fiber is the set of $v$ satisfying $Mv=0$. Over the integers, adding any such $v$ to one solution of $Mq=x$ gives another. For a meaning map the associated equivalence groups configurations with equal images.}
\sosterm{g-lex}{Lexicographic order, $<_{\rm lex}$}{Compare two coordinate vectors at their first differing position, from left to right. The vector with the smaller coordinate there comes first. Natural OCA uses this comparison only after sizes tie, giving every configuration finitely many predecessors. The total order is denoted $\prec$.}
\sosterm{g-matrix}{Multiplication matrix, $M_\varepsilon(y)$}{For a fixed divisor $y$, column $i$ of $M_{\rm R}(y)$ is $e_i\star y$. Column $i$ of $M_{\rm L}(y)$ is $y\star e_i$. With either choice, $M_\varepsilon(y)q=x$ is the quotient equation. The column sum $c_i^\varepsilon(y)$ counts output occurrences contributed by one unit at $i$. Every natural solution satisfies $\sum_i c_i^\varepsilon(y)q_i=\nu(x)$. \linesref{T}{377}{421}.}
\sosterm{g-multiset}{Multiset, multiplicity and occurrence}{An occurrence is one instance of a primitive type. Its type's multiplicity is the number of instances of that type. A multiset records each multiplicity, so changing the order in which occurrences are listed preserves configuration identity.}
\sosterm{g-name}{Name}{The exact original text of a structure-substrate. Registered names are nonempty. Name equality compares the sequence of Unicode scalar values, including every space, in order. Looking up a name and encoding that name as an integer are independent operations.}
\sosterm{g-namespace}{Namespace and profile, $\lambda$}{A namespace fixes the configuration domain, its equality criterion and its indexing convention, including alphabet order. A profile also supplies the composition rules and any admissibility constraints. The record stores the profile's text key $\lambda$. A bare address can denote different vectors in different namespaces.}
\sosterm{g-neutral-toggle}{Neutral and toggle}{Two primitive types in the running instance, ordered neutral first. Neutral is the identity, and toggle composed with toggle gives neutral. For vectors $(p,q),(r,s)$ the product is $(pr+qs,ps+qr)$. The mixed vector $(1,1)$ has natural rank $4$. \linesref{T}{889}{934}.}

\sosdictgroup{g-o}{O--Q}
\sosterm{g-oca}{Object Counting Algebra, $\mathfrak O(A,T,\mathscr S,\sigma)$}{The specification collecting configurations and their operations, address encoding and decoding, the transported operations, semantic interpretation and the typed evaluator. The parameters state which primitive types exist, what their pairs produce, and how results are interpreted. \linesref{T}{48}{66}.}
\sosterm{g-observer}{Observer origin and view, $o,V_o$}{An origin names a declared way of observing a record. Its view keeps the name and shape and applies a meaning map $\sigma_o$ to each cell. A type-counting view can return $(1,1)$ where a size-only view returns $2$. The record survives either observation.}
\sosterm{g-operation}{Operation and operand}{An operation is a rule taking inputs, called operands, to a result when defined. The symbols $x,y,z,q$ usually denote configurations and $a,b$ their integer addresses. An \emph{ordered} pair remembers which operand is first. The side of a divisor matters when multiplication can be noncommutative.}
\sosterm{g-primitive}{Primitive type and primitive rule}{A type $o_i$ is one element of the alphabet. Its singleton configuration is $e_i$. The rule $\mu_T$ assigns a whole output configuration to each ordered pair of types. It may produce zero, one or many occurrences. The finite monoid specialization always produces one.}
\sosterm{g-quotient}{Quotient set, $Q_{\rm R},Q_{\rm L}$}{A quotient is an admissible missing factor. $Q_{\rm R}(x,y)$ collects all $q$ with $q\star y=x$, whereas $Q_{\rm L}(x,y)$ collects all $q$ with $y\star q=x$. The subscript locates the \emph{known} divisor $y$. For natural mixed neutral--toggle divided by itself, both encoded sets equal $\{1,2\}$. A symbolic family denotes a set by a rule generating exactly its members.}

\sosdictgroup{g-r}{R--S}
\sosterm{g-record}{Record or structure-substrate, $R$}{The tuple $R=(s,\lambda,h,B)$ holds an exact name, profile key, shape and body. The body maps every valid coordinate to a configuration. A record with a zero extent has no cells. A record with one cell containing $\mathbf0$ has that cell and its empty configuration.}
\sosterm{g-registry}{Registry, $G,L,P$}{The outer array and its name index. The grid $G$ contains the permitted registry positions. The partial injective map $L$ sends a name to an occupied position. The partial array $P$ stores the record there. The conditions $\operatorname{im}L=\operatorname{dom}P$ and $P(L(s)).\mathrm{name}=s$ give exactly one name per occupied position. See \hyperlink{impl-2}{Implementation~2}.}
\sosterm{g-relabel}{Relabeling, $\pi$}{A bijection on the address space changes the encoder to $\rho_\pi=\pi\circ\rho$ and the decoder to $\delta_\pi=\delta\circ\pi^{-1}$. Transport every operation, constant, solution set and interpretation with this change. Then corresponding operands have corresponding results. The symbol $\cong_\pi$ denotes this isomorphism. \linesref{T}{640}{668}.}
\sosterm{g-residual}{Residual, $x\ominus y$ and $\mathcal R(x,y)$}{The exact additive remainder is the unique vector $z$ with $z\oplus y=x$, if one exists. In the natural profile it exists exactly when $y\preccurlyeq x$, and then $z_i=x_i-y_i$. The residual \emph{set} $\mathcal R(x,y)$ contains that vector when defined and is empty otherwise.}
\sosterm{g-restricted}{Restricted domain, $D$}{A declared subset of $\mathcal C$ containing zero. For $\omega_D$, inputs and result must belong to $D$ and the ambient operation must be defined. $Q_{\varepsilon,D}=Q_\varepsilon\cap D$ on admitted operands. The rank $\rho_D$ counts predecessors in $D$ alone. $I_D=\rho_D[D]$ and $\delta_D=\rho_D^{-1}$. If $|D|=k$ is finite, its addresses are $0,\ldots,k-1$. \linesref{T}{753}{810}.}
\sosterm{g-ring}{Ring}{A semiring in which every element has an additive inverse. For the signed algebra $\mathsf A_T$ and existence of a two-sided unit suffice to give a unital ring, written $\mathsf{Ring}_1$ in T. Commutativity of multiplication is a separate property.}
\sosterm{g-semantics}{Semantic interpretation and compatibility}{The target $\mathscr S=(S,\oplus_S,\star_S,0_S)$ gives possible meanings, their operations and designated zero. The map $\sigma:\mathcal C\to S$ assigns meanings. $\sigma_I=\sigma\circ\delta$ interprets addresses. Predicate $\mathsf H_\sigma$ requires preservation of zero, sums and products. $\mathcal R_S,Q_{\rm R}^{S},Q_{\rm L}^{S}$ solve the corresponding reconstruction equations in $S$. Under $\mathsf H_\sigma$ every interpreted configuration solution is a semantic solution (soundness). The reverse inclusion would require every such semantic solution to be the image of an admitted configuration solution (completeness). \linesref{T}{539}{599}.}
\sosterm{g-semiring}{Semiring}{A unital semiring has associative, commutative addition with zero, associative multiplication with a unit, distributivity on both sides, and multiplicatively absorbing zero. $\mathsf{Semiring}_1$ denotes this class. $\mathsf{CSemiring}_1$ adds commutativity of multiplication. \linesref{T}{348}{366}.}
\sosterm{g-shape}{Shape, extent and coordinate}{The shape $h=(h_1,\ldots,h_n)$ gives the number of positions along each of the $n$ axes. A coordinate $k$ satisfies $0\le k_j<h_j$ on every axis. There are $\prod_j h_j$ cells. In shape $(2,3)$ the last cell in row-major order has coordinate $(1,2)$. Here $n$ counts axes, while $d$ counts primitive types in each cell.}
\sosterm{g-shell}{Shell and combination count}{The natural size shell $\mathcal C_n$ contains the vectors of size $n$. Its size $\binom{n+d-1}{d-1}$ counts arrangements of $n$ identical occurrences separated into $d$ coordinates by $d-1$ dividers. For $0\le k\le m$, $\binom{m}{k}$ counts selections of $k$ positions among $m$. T uses zero when $m<k$ and an empty sum equals zero.}
\sosterm{g-signed}{Signed configuration, $\mathcal C^\pm$}{An integer vector $z\in\mathbb Z^d$ has possibly negative coefficients. Subtraction is total. Composition uses the same bilinear formula and division admits all integer solutions. $\|z\|_1=\sum_i|z_i|$ is its absolute size. The order $\prec_\pm$ uses this size and then lexicographic order to give the new encoder and decoder $\rho_\pm,\delta_\pm$.}
\sosterm{g-snapshot}{Snapshot}{A fixed valid registry state used throughout one read or trace. An immutable snapshot is held constant after publication. A later registration can publish a new snapshot. A caller using the old one continues to see its original name index and body values.}
\sosterm{g-sos}{Structure-substrate query, \texttt{structure\_substrate}}{The adopted read checks the arguments, resolves the name in the snapshot, then returns the pair containing the same name and its completion-derived learning flag. Reading a view or evaluating an algebra operation is a separate request. The original draft's entry point supplies completion $0$.}

\sosdictgroup{g-t}{T--Z}
\sosterm{g-transport}{Transport and hatted notation}{A hat marks an operation on addresses. For example, $a\widehat\star b=\rho(\delta(a)\star\delta(b))$. A quotient set is encoded member by member, $\widehat Q_\varepsilon(a,b)=\rho[Q_\varepsilon(\delta(a),\delta(b))]$. The rule works for empty and infinite sets as well. \linesref{T}{433}{476}.}
\sosterm{g-utf8}{Unicode scalar and UTF-8}{A scalar is an integer from hexadecimal \texttt{0000} through \texttt{10FFFF} excluding the surrogate range \texttt{D800}--\texttt{DFFF}. UTF-8 uses the shortest suitable byte pattern. The first byte has form \texttt{0xxxxxxx}, \texttt{110xxxxx}, \texttt{1110xxxx}, or \texttt{11110xxx} for lengths one through four. The remaining bytes have form \texttt{10xxxxxx}. Put the scalar's bits into the $x$ positions from most to least significant, left-padding with zeros. A decoder checks length, continuation bytes, shortest form and scalar range.}
\sosterm{g-view}{View}{The result of applying a particular observer's declared projection to a record. Views encode which detail is visible. Recovering the original cell from a projected value requires that its projection distinguish the relevant configurations. See \hyperlink{g-observer}{observer origin}.}
\sosterm{g-zero-column}{Zero column}{A column of $M_\varepsilon(y)$ whose entries all vanish. Over natural coefficients this is equivalent to $c_i^\varepsilon(y)=0$. If a solution exists, its coordinate $i$ can then be increased arbitrarily, giving countably many solutions. With no zero column, each $q_i$ is bounded above by $\lfloor\nu(x)/c_i^\varepsilon(y)\rfloor$.}

\manualtarget{g-notation}\subsection*{Mathematical signs}
\begin{description}
\item[Numbers and indices.] $\mathbb N_0=\{0,1,2,\ldots\}$ and $\mathbb Z$ includes all integers. $[d]=\{1,\ldots,d\}$. The letters $i,j,k,\ell,m$ can index coordinate positions. $\mathbf1_{\{i=k\}}$ is $1$ when the condition holds, $0$ otherwise. $\lfloor t\rfloor$ is the greatest integer at most $t$. The symbol $r\mid s$ means $s=rk$ for some integer $k$, and $r\nmid s$ means no such integer exists.
\item[Sets.] $\in$ means membership, $\subseteq$ inclusion, $\cap$ intersection, $\varnothing$ the empty set, and $\sqcup$ a disjoint union. $\mathcal P(U)$ contains all subsets of $U$. $\{q:F(q)\}$ denotes every $q$ satisfying condition $F$. $\aleph_0$ is the cardinality of $\mathbb N_0$. A set of this size can be put in bijection with the nonnegative integers.
\item[Maps and products.] $f:X\to Y$ gives one value in $Y$ for each member of $X$. The arrow $\rightharpoonup$ allows partiality. $f[U]=\{f(u):u\in U\}$. $g\circ f$ applies $f$ first. $\operatorname{id}_X$ returns its input. $X^d$ is the set of $d$-tuples over $X$. $\prod_j X_j$ is the set of tuples choosing one member from each $X_j$. $\sum$ adds terms and $\prod$ multiplies numeric factors.
\item[Assertions.] $:=$ introduces a definition, $\forall$ means for every, $\exists$ means there exists, and $\exists!$ means there is exactly one. $\land$ joins simultaneous requirements, $\lor$ allows either, $\Rightarrow$ denotes implication, and $\iff$ denotes implication in both directions. $\cong$ denotes an invertible correspondence preserving the stated operations or structure.
\end{description}

